\chapter{Resultados}
\label{cha:resultados}

\section{Resultados da busca bibliográfica}

A busca bibliográfica realizada em 29 de setembro de 2026 confirmou a existência de um conjunto relevante e relativamente recente de revisões sobre o metaverso, especialmente em torno de quatro eixos: definição conceitual, arquitetura, tecnologias e aplicações.

A literatura de Ritterbusch e Teichmann concentra-se na definição do conceito e identificou 28 definições ou descrições após o processo de seleção e buscas de citações. Tukur et al. concentraram-se nas técnicas, tecnologias e aplicações e chegaram a 30 estudos primários em sua revisão de escopo. Mais recentemente, Boztas et al. analisaram 103 estudos, dos quais 33 continham definições explícitas e 25 apresentavam modelos arquiteturais. Esses resultados mostram que a discussão acadêmica já ultrapassou a fase de simplesmente perguntar o que é o metaverso e passou a discutir também como seus ambientes podem ser estruturados e implementados \cite{ritterbusch2023,tukur2024,boztas2025}.

Zhou, Chen e Jin acrescentam uma perspectiva sociotécnica ao relacionar definições, características técnicas e comportamento dos usuários. Bénaben, Congès e Fertier propõem uma definição apoiada na evolução das tecnologias imersivas e em um framework de propriedades. Enamorado-Díaz et al. analisam aplicações, tecnologias e desafios de engenharia de software \cite{zhou2024,benaben2025,enamorado2025}.

\section{Síntese temática}

A comparação dos trabalhos permite organizar a literatura em três níveis:

\begin{enumerate}
    \item \textbf{Nível conceitual:} o que caracteriza um metaverso, incluindo ambientes digitais compartilhados, interação, identidade, imersão e persistência;
    \item \textbf{Nível tecnológico:} tecnologias e infraestruturas que podem viabilizar essas características, incluindo XR, inteligência artificial, computação distribuída e tecnologias descentralizadas;
    \item \textbf{Nível aplicacional:} usos em educação, saúde, manufatura, negócios, entretenimento e outros setores.
\end{enumerate}

Essa organização sustenta a principal decisão conceitual do trabalho: tecnologias não devem ser automaticamente convertidas em requisitos da definição. A realidade virtual, por exemplo, pode intensificar a imersão, mas não precisa ser tratada como condição suficiente ou necessária para toda manifestação do metaverso.

\section{Comparação com a definição do projeto}

A definição originalmente adotada no projeto apresenta forte convergência com a literatura ao considerar ambientes tridimensionais compartilhados, identidades digitais, interação, imersão e persistência. A interoperabilidade também aparece como dimensão relevante em trabalhos recentes, embora permaneça associada a desafios técnicos e de padronização.

A principal contribuição da revisão é, portanto, reorganizar a definição em camadas. A camada conceitual descreve o fenômeno; a camada tecnológica explica os mecanismos que podem sustentá-lo; e a camada aplicacional descreve os contextos nos quais o ecossistema pode ser utilizado.

\section{Limitação dos resultados quantitativos}

Os números publicados por outras revisões foram utilizados apenas como referência metodológica. O presente trabalho não atribui a si esses quantitativos.

O fluxograma PRISMA numérico desta pesquisa permanece dependente da execução das strings diretamente nas interfaces das bases selecionadas e do registro dos resultados retornados em cada uma delas. Essa limitação é explicitada para preservar a validade do artigo e evitar que resultados de estudos anteriores sejam apresentados como se fossem dados coletados nesta revisão.
